% ============================================================================
% DT-GSK equation registry (E1a-E12) -- LaTeX rendering (Phase 7, task C.1)
% Canonical source: papers/build_prompt_phases/phase_03/equation_registry.csv
% Symbols: papers/build_prompt_phases/phase_03/notation_table.md (tab:notation)
% MDPI-compatible: standard amsmath only. Zero empirical values.
% Each block is a faithful transcription of one registry row; the eq_id and
% code anchor are kept as comments. Do NOT edit equations here without a
% matching change-request against the frozen registry.
% ============================================================================

% ---- E1a (inherited) --- gained_shared_junior.py:56 ------------------------
% Junior knowledge-source indices (rank-neighbors; population sorted by
% fitness, r_i = rank of individual i).
\begin{equation}\label{eq:junior-idx}
R_1 = \operatorname{rank}(r_i - 1),\qquad
R_2 = \operatorname{rank}(r_i + 1),\qquad
R_3 \sim \mathcal{U}\!\left(P \setminus \{i, R_1, R_2\}\right)
\end{equation}
At the rank boundaries the neighbour pair is taken from the two nearest interior ranks: $(R_1,R_2)=(\operatorname{rank}(1),\operatorname{rank}(2))$ for the best individual ($r_i=0$) and $(R_1,R_2)=(\operatorname{rank}(NP{-}3),\operatorname{rank}(NP{-}2))$ for the worst ($r_i=NP{-}1$), as in the reference GSK implementation~\cite{mohamed2020gaining}.

% ---- E1b (inherited) --- gained_shared_senior.py:59 ------------------------
% Senior knowledge-source indices (top/mid/worst groups of size round(p NP)).
\begin{equation}\label{eq:senior-idx}
R_1 \sim \mathcal{U}(P_{\text{best}}),\quad
R_2 \sim \mathcal{U}(P_{\text{mid}}),\quad
R_3 \sim \mathcal{U}(P_{\text{worst}}),\qquad
|P_{\text{best}}| = |P_{\text{worst}}| = \operatorname{round}(p \cdot NP)
\end{equation}
The fraction $p$ is population-wide below $D=50$. At $D\ge50$ the frozen profile enables a conditional two-rate partition: while the trailing ten-generation mean acceptance is at most $0.10$, the worse-ranked half of the population draws its senior sources with $p=0.10$ and the better-ranked half keeps $p=0.05$; whenever acceptance is above that floor the single rate $p=0.05$ is used for the whole population.

% ---- E2 (inherited) --- _dt_core.py (Kexp path) ---------------------------
% Junior dimension schedule.
\begin{equation}\label{eq:kexp-schedule}
D_{\text{jun}} = \left\lceil D\,(1-x)^{K_{\exp}} \right\rceil,
\qquad p_{\text{jun}} = \frac{D_{\text{jun}}}{D},
\qquad x = \frac{t}{\text{MaxFES}}
\end{equation}
Each coordinate joins the junior phase with marginal probability $p_{\text{jun}}$, so $D_{\text{jun}}$ is the \emph{expected}, not the realized, junior-coordinate count. The assignment is an independent per-coordinate Bernoulli draw only on the rows that take the scalar mask; on a linkage-blockwise row (Eq.~(\ref{eq:kr-mask})) a single draw is shared across a whole block, so coordinates within a block are perfectly correlated and the realized count is more dispersed than the independent case.

% ---- E3 (inherited) --- _numba_accel.py:407-414 (gsk kernel) ---------------
% Gaining-sharing update. Junior and senior use DIFFERENT donor arrangements,
% exactly as in the original GSK: the JUNIOR phase differences the two
% rank-neighbours R1,R2 (Eq. junior-idx) and fitness-compares R3 to x_i; the
% SENIOR phase differences best(R1) and worst(R3) (Eq. senior-idx) and
% fitness-compares the MIDDLE source R2 to x_i. Because the two phases compare
% DIFFERENT donors, the sign is per-phase: s_J (junior, vs R3) and s_S (senior,
% vs R2), each in {+1,-1} (+1 when the compared source is fitter than x_i).
% R-01: the two signs are separated and DEFINED inside this SAME numbered
% display, and the convention is restated below in rendered text (it used to
% live only in this comment, so it reached neither the PDF nor the DOCX).
\begin{equation}\label{eq:gsk-update}
\begin{aligned}
\text{junior:}\quad u_i &= x_i + KF\left[\, (x_{R_1} - x_{R_2}) + s_J\,(x_{R_3} - x_i) \,\right],\\[2pt]
\text{senior:}\quad u_i &= x_i + KF\left[\, (x_{R_1} - x_{R_3}) + s_S\,(x_{R_2} - x_i) \,\right],\\[4pt]
s_J &= +1 \ \text{ if } f(x_i) > f(x_{R_3}), \ \text{ else } -1,\\[2pt]
s_S &= +1 \ \text{ if } f(x_i) > f(x_{R_2}), \ \text{ else } -1
\end{aligned}
\end{equation}

Here $s_J,s_S\in\{+1,-1\}$ are the signs of two \emph{different} fitness
comparisons: the junior phase compares the random source $x_{R_3}$ to $x_i$,
whereas the senior phase compares the middle source $x_{R_2}$ to $x_i$.
Both follow the same convention---$+1$ when the compared source is fitter
than $x_i$, so that the step is oriented towards it, and $-1$ otherwise
(an exact tie takes $-1$).

% ---- E4 (modified) --- _make_linkage_groups:918 / mask build --------------
% Crossover mask: per-coordinate Bernoulli(KR), OR a linkage-block mask that
% copies a whole block together. The random linkage blocks are permutation
% chunks -- arbitrary coordinate subsets, NOT contiguous index ranges; at
% D>=50 the SGSM graph's learned components (also arbitrary index sets) can
% supply the blocks.
\begin{equation}\label{eq:kr-mask}
v_{ij} =
\begin{cases}
u_{ij} & \text{if } m_{ij} = 1\\[2pt]
x_{ij} & \text{otherwise}
\end{cases}
\qquad
m_{ij} \sim
\begin{cases}
\operatorname{Bernoulli}(KR) & \text{per-coordinate mask}\\[2pt]
\text{shared across a linkage block } B\ni j & \text{linkage-blockwise mode}
\end{cases}
\end{equation}
Here $KR$ is not the single frozen constant of Table~\ref{tab:notation} but the rate carried by the ACE arm that individual $i$ drew this generation, so it varies per individual and per generation over the arm pool. It is additionally floored at $\min(1,\ \kappa_{\dim}/D)$ with $\kappa_{\dim}=0/20/40/40$ by tier, which raises the low-$KR$ arms above their nominal rates at every $D\ge20$. Finally, a force-nonzero guarantee applies after mask assembly in both modes: if a row's mask would leave the trial identical to its parent, one coordinate is set at random so that every trial differs from its parent in at least one coordinate.

% ---- E5 (modified) --- _dt_core.py:790 / _numba_accel.py:909 --------------
% Nonlinear population reduction (tier-floored; rounded half-up).
\begin{equation}\label{eq:nlpsr}
NP(x) = \operatorname{round}_{\uparrow\frac12}\!\left(
NP_{\text{init}} + \left(N_{\min} - NP_{\text{init}}\right) x^{(1-x)}
\right),
\qquad x = \frac{t}{\text{MaxFES}}
\end{equation}

% ---- E6 (modified) --- _ace_update_probs:1203 / _ace_project_probs:694 -----
% ACE probability update, transcribed from the frozen operator. Paper convention
% d_m = -omega_m and S = -s, where omega/s are the code's variables. Three
% branches: net improvement (S>0) credits the FULL arm pool d/S; a non-improving
% generation (S<0) with at least one improving arm credits that subset only;
% S=0, non-finite S, or no improving arm holds pi (projection only). The target
% is projected BEFORE the EMA mix and the mixture projected again; the projection
% is Euclidean onto the floored simplex, not a floor-and-renormalize.
\begin{equation}\label{eq:ace-update}
\begin{gathered}
d_m = \!\!\sum_{i:\,s_i = m}\!\! \bigl(f(x_i) - f(v_i)\bigr),
\qquad
S = \sum_{m} d_m,
\qquad
\mathcal{P}(z) = \operatorname*{arg\,min}_{\substack{y_m \ge \pi_{\min}\\ \sum_{m} y_m = 1}} \lVert y - z \rVert_2,
\\[4pt]
\tau =
\begin{cases}
d \,/\, S, & S > 0,\\[2pt]
\bigl(\max(d_m,0)\bigr)_m \Big/ \sum_{m'}\max(d_{m'},0), & S < 0 \ \text{ and } \ \max_m d_m > 0,
\end{cases}
\\[4pt]
\pi \leftarrow
\begin{cases}
\mathcal{P}\bigl((1-c)\,\pi + c\,\mathcal{P}(\tau)\bigr), & \tau \text{ defined},\\[2pt]
\mathcal{P}(\pi), & \text{otherwise.}
\end{cases}
\end{gathered}
\end{equation}

% ---- E7 (inherited) --- bound_constraint.py:1 ------------------------------
% L-SHADE midpoint bound repair.
\begin{equation}\label{eq:midpoint-repair}
v_{ij} \leftarrow
\begin{cases}
\dfrac{x_{ij} + \ell_j}{2} & \text{if } v_{ij} < \ell_j\\[6pt]
\dfrac{x_{ij} + u_j}{2} & \text{if } v_{ij} > u_j
\end{cases}
\end{equation}

% ---- E8 (inherited) --- dt_gsk_optimize accept step (_dt_core.py:3325) ----
% Greedy selection: the trial replaces the parent when it is no worse (ties
% accepted, <=), as in the inherited GSK operator. At D>=100 a conditional
% late-budget acceptance clip can tighten this to a strict relative decrease
% over a stall streak (the high-dimensional late-accept safeguard).
\begin{equation}\label{eq:greedy-accept}
(x_i, f_i) \leftarrow \left(v_i,\, f(v_i)\right)
\quad\text{iff}\quad f(v_i) \le f(x_i)
\end{equation}
The tie-accepting form above is the base rule and is what runs at $D<100$. At $D\ge100$ the late-acceptance clip can replace it, late in the budget and after a stall streak, with the strict test $f(v_i) < f(x_i) - \epsilon\lvert f(x_i)\rvert$ at $\epsilon=10^{-3}$, which rejects ties; the clip's own gates are given in Supplementary Section~S5.

% ---- E9 (modified) --- _dt_core.py:3874 (cauchy rescue) -------------------
% Budget-safe escape: Cauchy rescue of the worst round(r_c NP) UNFROZEN rows
% (the top n_freeze elites are protected). Per-coordinate scale
% gamma = r_s (u-l)/sqrt(D) (r_s=bse_cauchy_scale); draws are bound-repaired,
% charged to the budget, and kept only if not worse (<=).
\begin{equation}\label{eq:bse-cauchy}
x_i \leftarrow \operatorname{repair}\!\bigl(x_i + \gamma\,\mathcal{C}(0,1)\bigr),
\quad \gamma = r_{s}\,\frac{u-\ell}{\sqrt{D}},
\quad\text{worst } \operatorname{round}(r_{c}\, NP)\text{ unfrozen rows, kept iff } \le
\end{equation}

% ---- E10 (original) --- _dt_subsystems/interaction_graph.py:222 -----------
% SGSM interaction memory: four decaying accumulators (magnitude G_abs, signed
% G_signed, per-dimension activity g_act, and co-occurrence support C) EMA-updated
% from STRICTLY-IMPROVING accepted moves only (ties and DE-arm moves excluded;
% local-search moves weighted 0.25). Each displacement is L1-normalized over its
% active coordinates and weighted by its improvement w_i; the graph acts only
% when conf(G_abs) reaches kappa_min. The full per-sample normalization, source
% weighting, and confidence functional are stated normatively in Supplement S5.
% M-010: this equation is the UPDATE rule only. Consumption (whether a learned
% block or eigenbasis may be USED) is separately gated on conf(G_abs) >= kappa_min
% and is stated in the prose, not here -- updating and consuming are distinct.
\begin{equation}\label{eq:sgsm-graph}
\begin{aligned}
G_{\bullet} &\leftarrow \lambda\,G_{\bullet} +
\eta \sum_{i\,\in\,\mathcal{I}^{+}} w_i\, \phi_{\bullet}(\hat{\delta}_i)\, \phi_{\bullet}(\hat{\delta}_i)^{\!\top},
\qquad \hat{\delta}_i = \frac{\Delta x_i}{\lVert \Delta x_i\rVert_1},\\[2pt]
&\bigl(\bullet\in\{\mathrm{abs},\mathrm{signed}\};\ \phi_{\mathrm{abs}}{=}|\cdot|,\ \phi_{\mathrm{signed}}{=}\mathrm{id}\bigr),
\end{aligned}
\end{equation}

% ---- E11 (original) --- _final_polish_basis:1882 / _final_polish_compass:1907
% Eigenframe compass polish: one-shot, RNG-free, strict budget accounting;
% V = eigenbasis of the SGSM signed matrix (coordinate axes if no signal).
\begin{equation}\label{eq:eigen-polish}
x^{*} \leftarrow \operatorname{CompassSearch}\!\left(x^{*},\,
V = \operatorname{eig}(G_{\mathrm{signed}})\right)
\end{equation}

% ---- E12 (original) --- _dt_rng.py:1 --------------------------------------
% RNG substream seeding: 13 append-only, prefix-locked substreams; stream 0
% is the run seed verbatim; child seeds are a counter offset mod M, then
% each stream is an independent Threefry generator from its child seed.
\begin{equation}\label{eq:rng-substreams}
\begin{aligned}
\text{stream}_0 &= \operatorname{threefry}(\text{seed}),
\qquad
\text{stream}_k = \operatorname{threefry}(s_k),\\
s_k &= (\text{seed} + 1{,}000{,}003\,(k{+}1)) \bmod M_{\mathrm{seed}} + 1,
\quad k \in \{1,\dots,12\}
\end{aligned}
\end{equation}
where $M_{\mathrm{seed}} = 2\,147\,483\,646$ is the child-seed modulus
(Supplementary Material, Section~S5).
